\documentclass[10pt,twocolumn]{article}

% -----------------------------------------------------------------------
% PACKAGES
% -----------------------------------------------------------------------
\usepackage[T1]{fontenc}
\usepackage[utf8]{inputenc}
\usepackage{lmodern}
\usepackage{microtype}
\usepackage{geometry}
\geometry{
    a4paper,
    top=1.8cm, bottom=2.2cm,
    left=1.6cm, right=1.6cm,
    columnsep=0.6cm
}

\usepackage{amsmath,amssymb,amsfonts}
\usepackage{mathtools}
\usepackage{bm}
\usepackage{graphicx}
\usepackage{float}
\usepackage{subcaption}
\usepackage{booktabs}
\usepackage{array}
\usepackage{xcolor}
\usepackage{url}
\usepackage{hyperref}
\hypersetup{
    colorlinks=true,
    linkcolor=blue!60!black,
    citecolor=blue!60!black,
    urlcolor=blue!60!black
}
\usepackage[numbers,sort&compress]{natbib}
\usepackage{algorithm}
\usepackage{algpseudocode}
\usepackage{tikz}
\usetikzlibrary{shapes.geometric, arrows.meta, positioning, fit, backgrounds}
\usepackage{pgfplots}
\pgfplotsset{compat=1.18}
\usepackage{enumitem}
\usepackage{caption}
\captionsetup{font=small, labelfont=bf}
\usepackage{fancyhdr}
\pagestyle{fancy}
\fancyhf{}
\fancyhead[L]{\small\textit{PureDL CT Reconstruction --- Workshop Submission 2025}}
\fancyhead[R]{\small\thepage}
\renewcommand{\headrulewidth}{0.4pt}

% -----------------------------------------------------------------------
% CUSTOM COMMANDS
% -----------------------------------------------------------------------
\newcommand{\todo}[1]{\textcolor{red}{\textbf{[TODO: #1]}}}
\newcommand{\radon}{\mathcal{R}}
\newcommand{\myvec}[1]{\bm{#1}}
\newcommand{\figplaceholder}[2]{%
    \fbox{\parbox[c][#2]{\linewidth-2\fboxsep}{%
        \centering\textit{\small #1}%
    }}%
}

% -----------------------------------------------------------------------
% TITLE
% -----------------------------------------------------------------------
\title{%
    \LARGE\bfseries
    Pure Deep Learning CT Reconstruction from Physically Simulated\\
    Cone-Beam Projections: A Three-Stage, Radon-Transform-Free Pipeline\\[4pt]
    \normalsize\mdseries
    \textit{Bypassing Analytical Inversion with a Learned Dual-Domain Architecture}
}

\author{%
    \begin{tabular}{c}
        \large [Author Name(s)] \\[2pt]
        \normalsize [Affiliation / University / Laboratory] \\
        \normalsize [Department, City, Country] \\
        \small \texttt{[email@institution.edu]}
    \end{tabular}
}

\date{July 2025}

% -----------------------------------------------------------------------
\begin{document}
% -----------------------------------------------------------------------

\maketitle
\thispagestyle{fancy}

% -----------------------------------------------------------------------
\begin{abstract}
% -----------------------------------------------------------------------
Computed Tomography (CT) reconstruction has historically relied on analytical
methods --- Filtered Back-Projection (FBP) and the Feldkamp-Davis-Kress (FDK)
algorithm --- whose accuracy degrades substantially under
sparse-view or limited-angle acquisition regimes, and which are
fundamentally tied to exact knowledge of the scanning geometry.
We present a \textit{Pure Deep Learning} (PureDL) pipeline that
\textbf{entirely replaces} the classical Radon-transform inversion step with
a three-stage, end-to-end trainable neural network.
The pipeline operates directly on raw sinograms and produces
high-quality 2D cross-sectional CT slices without invoking any
explicit back-projection operator during inference.
Our simulation infrastructure leverages \textbf{gVirtualXray}~(gVXR) to
generate physically realistic cone-beam projection data from
CAD-sourced STL meshes with full photon-counting noise modelling
(Poisson shot noise and Gaussian detector readout noise), providing a
principled, reproducible training environment.
The three-stage architecture consists of:
(i) a \textit{SinogramUNet} for sensor-domain denoising,
(ii) a dilated \textit{ConvolutionalDomainTransformer} that learns to
globally invert the Radon transform from data alone, and
(iii) an \textit{ImageUNet} for image-domain enhancement and
edge sharpening.
Training is supervised with a multi-term L1 loss that provides
simultaneous guidance in both the sinogram and image domains.
We detail the full software pipeline from geometry simulation to
network training and provide a comparative experimental framework
against the classical FDK baseline using ASTRA Toolbox.
[Results will be reported upon GPU execution completion.]
\end{abstract}

\textbf{Keywords:} computed tomography, deep learning, inverse problems,
sinogram processing, FDK reconstruction, gVXR simulation,
cone-beam CT, image reconstruction, U-Net, dilated convolution.

% -----------------------------------------------------------------------
\section{Introduction}
\label{sec:intro}
% -----------------------------------------------------------------------

Computed Tomography (CT) is an indispensable imaging modality in
medical diagnostics, materials science, and industrial non-destructive
testing (NDT).
Classical reconstruction algorithms, particularly Filtered
Back-Projection~\cite{kak2001} and its cone-beam adaptation by
Feldkamp, Davis, and Kress~\cite{feldkamp1984}, provide mathematically
exact reconstructions under ideal full-angle, low-noise acquisition.
However, these methods exhibit severe artefacts under three practically
relevant conditions: (1) sparse angular sampling (to reduce radiation
dose or scan time), (2) limited angular coverage (when physical
constraints prevent full $360^\circ$ rotation), and (3) high noise
levels arising from photon-starved acquisitions.

Beyond quality degradation, FDK reconstruction is inherently coupled
to the precise geometric calibration of the scanner.
Any miscalibration in the source-object distance (SOD) or
source-detector distance (SDD) propagates directly into the
reconstructed volume as ring artefacts, geometric distortions,
and blurring.
In real-world NDT scenarios involving industrial objects of
non-standard geometry --- such as turbine blades, castings, or
machined mechanical components --- this calibration requirement is
both expensive and error-prone.

The deep learning revolution has prompted the community to ask:
\textit{can a neural network learn to replace the physics-based
inversion step entirely?}
Early approaches~\cite{jin2017,chen2017} applied CNNs in the image
domain to post-process classical FDK outputs.
Subsequent works~\cite{han2016,lee2018} demonstrated that operating
directly on sinograms, or simultaneously in both domains~\cite{hu2020},
yields superior artefact suppression.
However, these methods all retain the classical back-projection as a
mandatory intermediate step.

\textbf{Our contribution} is a system that \textbf{eliminates the
classical back-projection entirely} at inference time.
We propose the \emph{PureDLPipeline} --- a three-stage, end-to-end
trained architecture that receives a raw, noisy sinogram and outputs
the final CT cross-section using only learned transformations.
The key challenge is that inverting the Radon transform requires
integrating information across the full angular range simultaneously,
which demands a truly global receptive field.
We address this by introducing a bottleneck sub-network with
hierarchically stacked dilated convolutions, whose receptive field
is analytically shown to grow exponentially with depth.

Furthermore, unlike prior work that trains on measured clinical or
industrial CT data (which suffers from limited ground-truth
availability), we leverage \textbf{gVirtualXray}~(gVXR)~\cite{vidal2019gvxr}
to generate a controlled, scalable dataset from open 3D meshes,
enabling precise control over degradation type and level.

\subsection{Summary of Contributions}

\begin{enumerate}[leftmargin=*, label=\textbf{\arabic*.}]
    \item \textbf{CDT: A resolution-independent, dilated domain transformer
    for Radon inversion.}
    Existing pure-DL reconstruction networks either use vanilla U-Nets
    that lack a principled global receptive field, or massive
    transformer architectures requiring prohibitive compute.
    We introduce the \textit{Convolutional Domain Transformer} (CDT):
    a compact encoder-dilated-bottleneck-decoder that (a)~explicitly
    solves the global receptive field problem through stacked
    dilations, and (b)~achieves resolution independence via an
    initial bilinear alignment, allowing the same trained model to
    operate on different detector pixel counts without retraining.
    This precise combination is not present in the prior literature.

    \item \textbf{First complete gVXR-to-PureDL closed-loop
    pipeline for industrial NDT objects.}
    While gVXR has been applied to segmentation and
    detection tasks~\cite{vidal2019gvxr}, no published work has
    demonstrated an end-to-end pipeline from STL mesh simulation
    through degradation generation to pure-DL reconstruction training
    and full-volume inference for industrial non-destructive testing
    geometries.
    Our pipeline handles centered and offset cone-beam geometries,
    multi-material simulation (Ti, Al), and systematically generates
    80+ training configurations from only 5 object meshes.

    \item \textbf{Stage-wise independently supervised training
    with three-term dual-domain L1 loss.}
    We apply independent image-domain supervision not only at the
    final output (standard in dual-domain methods) but also directly
    to the Stage~2 rough image.
    This provides an explicit gradient path into the CDT bottleneck
    that bypasses Stage~3, preventing gradient starvation at the
    most critical learned component.

    \item \textbf{Systematic diversity training across object
    geometries, materials, and degradation modes.}
    Our dataset design spans sparse-view ($\times$2, $\times$4),
    limited-angle (two angular windows), and noise-amplified
    acquisitions simultaneously across all training objects.
    This cross-condition exposure produces a single model capable
    of handling all three degradation types at inference,
    without per-condition retraining.
\end{enumerate}

\subsection{What We Do \emph{Not} Claim as Novel}
\label{sec:not_novel}

In the interest of rigour, we explicitly acknowledge the following:
\begin{itemize}[leftmargin=*, noitemsep]
    \item \textit{Pure sinogram-to-image DL mapping} is an active
    research direction; our contribution is the specific CDT
    architecture and the NDT application context, not the paradigm
    itself.
    \item \textit{Dual-domain training losses} have been used in
    DRONE~\cite{wu2021drone} and related work;
    our novelty is the stage-wise independent bottleneck supervision.
    \item \textit{Slice-by-slice 2D processing} is a deliberate
    engineering choice driven by GPU memory constraints.
    A 3D joint volumetric approach would improve cone-coupling
    accuracy and is the natural extension of this work.
    \item \textit{Dilated convolutions} are well-established;
    their specific application to global Radon inversion in a
    resolution-independent setting is our contribution.
\end{itemize}

% -----------------------------------------------------------------------
\section{Background and Related Work}
\label{sec:related}
% -----------------------------------------------------------------------

\subsection{The Radon Transform and its Inversion}

In 2D parallel-beam geometry, the measurement at angle $\theta$ and
detector position $t$ is the \textit{line integral} of the
attenuation coefficient $f(\myvec{x})$ along the ray
perpendicular to $\theta$:
\begin{equation}
    p(t, \theta) = \int_{-\infty}^{\infty}
        f(t\cos\theta - s\sin\theta,\; t\sin\theta + s\cos\theta)\;
    \mathrm{d}s
    \label{eq:radon}
\end{equation}
The full collection of projections $p(t, \theta)$ over
$\theta \in [0, \pi)$ is the \textit{sinogram}.
The analytical inverse is the Filtered Back-Projection:
\begin{equation}
    f(\myvec{x}) = \int_0^{\pi}
        \left(p(\cdot, \theta) * h\right)(x\cos\theta + y\sin\theta)
    \;\mathrm{d}\theta
    \label{eq:fbp}
\end{equation}
where $h$ is a ramp filter in the frequency domain.
In cone-beam geometry, FDK~\cite{feldkamp1984} approximates this
integral using cosine-weighted projections and a modified ramp filter.
Both methods require dense angular sampling ($>180^\circ + \text{fan angle}$)
and degrade predictably with fewer angles.

\subsection{Deep Learning for CT Reconstruction}

\textbf{Image-domain post-processing.}
Early works such as RED-CNN~\cite{chen2017} and FBPConvNet~\cite{jin2017}
apply CNNs to FDK output images to suppress streak artefacts from
sparse-view sampling.
These require a classical reconstruction as input, inheriting its
geometric constraints and calibration requirements.

\textbf{Sinogram-domain preprocessing.}
Approaches such as~\cite{lee2018} use CNNs to complete or denoise
sinograms \textit{before} FDK reconstruction, recovering missing
angular information in a learned manner.
However, analytical back-projection remains a mandatory downstream step,
so any geometric calibration error in the scanner is still
propagated into the final volume.

\textbf{Dual-domain methods.}
DRONE~\cite{wu2021drone} and related networks process data
simultaneously or sequentially in both the sinogram and image
domains.
They represent the current state-of-the-art in artefact suppression
but remain architecturally coupled to FBP/FDK: back-projection is
always invoked as an intermediate step, and the network is
consequently constrained to known, fixed scanner geometries.

\textbf{Physics-replacement and pure DL networks.}
iCTNet~\cite{liu2020} and similar works attempt to replace the
back-projection operator with learned transformations.
However, these approaches typically either:
(a) use vanilla U-Nets that lack a principled global receptive field
for angular integration, relying on depth alone to indirectly
expand the network's field of view;
(b) require explicit knowledge of the scanning geometry as a
conditioned input; or
(c) have only been demonstrated on idealised 2D parallel-beam
fantoms (Shepp-Logan) rather than realistic cone-beam datasets.

More recently, transformer-based architectures~\cite{han2016}
have shown promise for capturing global sinogram structure
but at a computational cost ($>$100\,M parameters, days of
training) that prohibits application to standard research
GPU budgets.

\textbf{Where our work sits.}
We occupy the gap between these extremes:
a \emph{compact} ($\approx$8.8\,M parameter) pure-DL architecture
designed specifically to solve the global receptive field problem
through dilated bottleneck convolutions, tested on realistic
cone-beam projections from industrial geometry rather than
medical phantoms.

\subsection{Simulation-Based Dataset Generation}

Training data for CT reconstruction networks requires paired
(degraded sinogram, ground-truth image) samples.
Options include: (a) retrospective clinical datasets
(limited availability, ethics-restricted),
(b) mathematical phantoms (Shepp-Logan, unrealistically smooth
and geometry-limited),
and (c) physics simulation from 3D models.
gVXR~\cite{vidal2019gvxr} provides GPU-accelerated ray-casting
X-ray simulation from arbitrary polygon meshes, enabling rapid
generation of large, diverse datasets with physically correct
photon statistics.

While gVXR has been applied to segmentation and
detection tasks in the NDT domain, its use as the exclusive
dataset engine for training a pure-DL \emph{reconstruction}
network has not been previously reported.
Table~\ref{tab:comparison} positions our work against
representative prior methods.

\begin{table}[H]
    \centering
    \caption{%
        Comparison with representative related methods.
        \checkmark~=~yes, $\times$~=~no.
        CDT = Convolutional Domain Transformer (ours).
    }
    \label{tab:comparison}
    \resizebox{\linewidth}{!}{%
    \begin{tabular}{lccccc}
        \toprule
        \textbf{Method} &
        \textbf{No FDK} &
        \textbf{Cone-beam} &
        \textbf{Industrial} &
        \textbf{Res.-indep.} &
        \textbf{Synthetic data} \\
        \midrule
        FBPConvNet~\cite{jin2017}   & $\times$ & $\times$ & $\times$ & $\times$ & $\times$ \\
        RED-CNN~\cite{chen2017}     & $\times$ & $\times$ & $\times$ & $\times$ & $\times$ \\
        DRONE~\cite{wu2021drone}    & $\times$ & $\times$ & $\times$ & $\times$ & $\times$ \\
        iCTNet~\cite{liu2020}       & \checkmark & $\times$ & $\times$ & $\times$ & \checkmark \\
        \textbf{Ours (CDT)}         & \checkmark & \checkmark & \checkmark & \checkmark & \checkmark \\
        \bottomrule
    \end{tabular}}
\end{table}

% -----------------------------------------------------------------------
\section{Physics Simulation and Dataset Generation}
\label{sec:data}
% -----------------------------------------------------------------------

\subsection{gVXR Simulation Framework}

We use \textbf{gVirtualXray}~(gVXR) version $\geq$\,3.0 to simulate
cone-beam X-ray projections.
gVXR operates on triangle meshes (STL format) and computes the
energy deposited on a virtual detector by casting rays through the
scene using Beer-Lambert attenuation:
\begin{equation}
    I = I_0 \cdot \exp\!\left(-\int \mu(\myvec{x})\,\mathrm{d}\ell\right)
    \label{eq:beer}
\end{equation}
where $\mu(\myvec{x})$ is the spatially-varying linear attenuation
coefficient determined by the material element assignment
(e.g., Titanium, Aluminium).

\paragraph{Scanner geometry.}
The virtual scanner replicates a Nikon XT\,H\,225\,ST setup:
\begin{itemize}[leftmargin=*,noitemsep]
    \item Source-Object Distance (SOD): 300\,mm
    \item Source-Detector Distance (SDD): 1100\,mm
    \item Detector: $1800 \times 1800$ pixels at 0.2\,mm/pixel
    \item Beam: 100\,kV monochromatic at 1000 photons/ray
\end{itemize}

\paragraph{Magnification and FOV.}
The magnification factor is $M = \mathrm{SDD}/\mathrm{SOD} = 3.67$.
The physical Field of View at the isocentre is:
\begin{equation}
    \mathrm{FOV}_{\mathrm{obj}} = \frac{N_{\mathrm{det}} \cdot p_{\mathrm{det}}}{M}
    = \frac{1800 \times 0.2}{3.67} \approx 98.1\;\mathrm{mm}
    \label{eq:fov}
\end{equation}
We use 90\% of this value as the safe FOV to avoid truncation artefacts.

\paragraph{Object preprocessing.}
Each STL mesh is automatically preprocessed by:
\begin{enumerate}[leftmargin=*,noitemsep]
    \item Centering the bounding box centroid at the isocentre.
    \item Aligning the longest axis to $Z$ (vertical detector axis).
    \item Scaling the mesh so its diagonal fits within the safe FOV.
\end{enumerate}
This ensures all objects are fully contained within the detector
without any a priori knowledge of their size.

\subsection{Noise Model}

The ideal transmission value $T = I/I_0 \in [0, 1]$ is corrupted by
two independent noise sources to simulate realistic detector physics:

\textbf{Poisson (photon shot) noise:}
\begin{equation}
    \tilde{N} \sim \mathrm{Poisson}\!\left(T \cdot I_0^{\mathrm{photon}}\right),
    \quad I_0^{\mathrm{photon}} = 50{,}000
    \label{eq:poisson}
\end{equation}

\textbf{Gaussian (electronic readout) noise:}
\begin{equation}
    \hat{N} = \tilde{N} + \mathcal{N}(0, \sigma_e^2),
    \quad \sigma_e = 10
    \label{eq:gaussian}
\end{equation}

The noisy attenuation measurement is then:
\begin{equation}
    p_{\mathrm{noisy}} = -\ln\!\left(\frac{\hat{N}}{I_0^{\mathrm{photon}}}\right)
    \label{eq:att}
\end{equation}

To prevent numerical instabilities, we clip $\hat{N} \geq 1$.
The resulting projections are scaled by per-projection maximum
attenuation and saved as 16-bit TIFF files along with the per-projection
scale factor array \texttt{sino\_scales.npy} for lossless reconstruction.

\subsection{Projection Acquisition Protocol}

We acquire $N_\theta = 360$ projections uniformly spaced over
$[0^\circ, 360^\circ)$ at $1^\circ$ angular step.
Each projection is an $1800 \times 1800$ pixel image.
From this full set, three degraded subsets are constructed to form
the training distribution:

\begin{enumerate}[leftmargin=*,noitemsep]
    \item \textbf{Sparse-view} ($\times 2$ and $\times 4$ undersampling):
    every 2nd or 4th projection is retained, yielding 180 or 90
    projections respectively.
    \item \textbf{Limited-angle}: only projections in $[0^\circ, 180^\circ]$
    or $[30^\circ, 210^\circ]$ are used.
    \item \textbf{High-noise}: full 360 projections with amplified
    Poisson noise ($I_0^{\mathrm{photon}} \in \{0.25 \times, 0.5 \times\}$ standard level).
\end{enumerate}

The paired ground-truth for each degraded sinogram is the FDK
reconstruction computed from the clean, full-angle sinogram by the
ASTRA Toolbox~\cite{astra2016}.

\begin{figure}[H]
    \centering
    \figplaceholder{%
        Figure 1: gVXR simulation pipeline.
        Left: STL mesh of an industrial component (whistle geometry).
        Centre: raw X-ray projection at $0^\circ$ and $90^\circ$.
        Right: full sinogram (Angles $\times$ Detectors) showing characteristic
        sinusoidal traces of high-density features.
        [RESULTS PENDING GPU EXECUTION]
    }{5cm}
    \caption{
        \textbf{Simulated Projection Data.}
        The gVXR engine produces physically realistic projections
        including Poisson photon noise and Gaussian detector readout
        noise. The sinogram (right) encodes the complete volumetric
        information in the measurement domain.
    }
    \label{fig:gvxr_data}
\end{figure}

% -----------------------------------------------------------------------
\section{Architecture: The PureDL Pipeline}
\label{sec:arch}
% -----------------------------------------------------------------------

The central hypothesis of this work is that a sufficiently expressive
neural network, trained on a large and diverse dataset of
(sinogram, image) pairs, can learn to approximate the inverse Radon
transform entirely from data, without requiring an explicit
analytical back-projection operator.

The \textbf{PureDLPipeline} decomposes this challenging mapping into
three specialised stages, each targeting a distinct sub-problem
(Figure~\ref{fig:arch_overview}).
Formally, let $\myvec{s} \in \mathbb{R}^{N_\theta \times N_d}$ be the
noisy input sinogram ($N_\theta$ angles, $N_d$ detector pixels).
The pipeline computes:
\begin{align}
    \hat{\myvec{s}}_{\text{clean}} &= \mathcal{F}_1(\myvec{s};\,\Theta_1)
        \label{eq:stage1} \\
    \myvec{r} &= \mathcal{F}_2(\hat{\myvec{s}}_{\text{clean}};\,\Theta_2)
        \in \mathbb{R}^{W \times H}
        \label{eq:stage2} \\
    \hat{\myvec{x}} &= \mathcal{F}_3(\myvec{r};\,\Theta_3)
        \label{eq:stage3}
\end{align}
where $\hat{\myvec{x}}$ is the final reconstructed CT slice and
$\Theta_1, \Theta_2, \Theta_3$ are the learnable parameters.

\begin{figure}[H]
    \centering
    \figplaceholder{%
        Figure 2: PureDLPipeline overview block diagram.
        Stage 1 (SinogramUNet): sinogram denoising.
        Stage 2 (ConvolutionalDomainTransformer): measurement-to-image domain mapping.
        Stage 3 (ImageUNet): image-domain sharpening and enhancement.
        Arrows showing data flow: noisy sinogram $\to$ clean sinogram
        $\to$ rough image $\to$ final CT slice.
        [DIAGRAM TO BE GENERATED]
    }{4.5cm}
    \caption{
        \textbf{PureDL Three-Stage Architecture.}
        The network processes information sequentially through a
        sinogram-domain denoiser, a learned domain transformer,
        and an image-domain refiner, without invoking any analytical
        back-projection operator at any stage.
    }
    \label{fig:arch_overview}
\end{figure}

\subsection{Stage 1: SinogramUNet --- Sensor Domain Rectification}
\label{sec:stage1}

\paragraph{Motivation.}
The Radon transform is a linear, invertible operator on smooth
functions.
However, the physical noise model (Section~\ref{sec:data}) breaks
linearity through the logarithmic transformation
($p = -\ln(I/I_0)$) and introduces signal-dependent Poisson noise
whose variance grows with attenuation.
If the corrupted sinogram is passed directly to the domain
transformer, residual noise patterns propagate into the final image
as structured artefacts that are significantly harder to remove
in the image domain.
Stage 1 performs a pre-conditioning of the sinogram in its native
measurement domain, where noise characteristics are more regular
and prior knowledge (smooth sinusoidal trace of each feature) can
be exploited.

\paragraph{Architecture.}
The SinogramUNet $\mathcal{F}_1$ is a lightweight encoder-decoder
with a single downsampling level (Figure~\ref{fig:stage1}):
\begin{align*}
    \myvec{f}_1 &= \mathrm{DoubleConv}_{1 \to C}(\myvec{s})
        \in \mathbb{R}^{C \times N_\theta \times N_d} \\
    \myvec{f}_2 &= \mathrm{DoubleConv}_{C \to 2C}(\mathrm{MaxPool}_{2}(\myvec{f}_1))
        \in \mathbb{R}^{2C \times \frac{N_\theta}{2} \times \frac{N_d}{2}} \\
    \myvec{f}_{\mathrm{up}} &= \mathrm{Bilinear}_{\times 2}(\myvec{f}_2)
        + \myvec{f}_1 \\
    \myvec{f}_3 &= \mathrm{DoubleConv}_{2C \to C}(\myvec{f}_{\mathrm{up}}) \\
    \hat{\myvec{s}} &= \mathrm{Conv}_{C \to 1}(\myvec{f}_3) + \myvec{s}
\end{align*}
where $C = 32$ and \texttt{DoubleConv} is a standard block:
\begin{equation}
    \mathrm{DoubleConv}_{a \to b}(\cdot) =
    \mathrm{ReLU}\!\circ\!\mathrm{BN}\!\circ\!\mathrm{Conv}_{3\times3}^{a \to b}
    \circ
    \mathrm{ReLU}\!\circ\!\mathrm{BN}\!\circ\!\mathrm{Conv}_{3\times3}^{a \to b}
    \label{eq:doubleconv}
\end{equation}

\paragraph{Residual connection.}
The final output $\hat{\myvec{s}} = \mathrm{Conv}(\myvec{f}_3) + \myvec{s}$
employs a global skip connection from input to output.
This design choice is deliberate: it ensures that in the
worst case (when all learned weights are zero), the network passes
its input through unchanged, guaranteeing that gradient flow is
stable at initialisation and that the network focuses on learning
the \textit{residual} correction rather than reconstructing the
entire signal from scratch.

\begin{figure}[H]
    \centering
    \figplaceholder{%
        Figure 3: SinogramUNet detailed architecture diagram.
        Show: Input sinogram (360$\times$1800 example) $\to$
        DoubleConv encoder $\to$ MaxPool $\to$ DoubleConv bottleneck
        $\to$ Bilinear upsample $\to$ skip-add $\to$ DoubleConv decoder
        $\to$ 1$\times$1 conv $\to$ residual add $\to$ clean sinogram output.
        [DIAGRAM PENDING]
    }{4cm}
    \caption{
        \textbf{Stage 1: SinogramUNet.}
        A single-resolution U-Net with a residual bypass.
        The network predicts a noise correction map in sinogram space,
        which is added back to the corrupted input.
    }
    \label{fig:stage1}
\end{figure}

\subsection{Stage 2: ConvolutionalDomainTransformer --- Physics Replacement}
\label{sec:stage2}

This is the \textbf{critical} and most novel component of the
pipeline.
It must learn to invert the Radon transform from data, an operation
that classically requires integrating information over the entire
angular range simultaneously.

\paragraph{The global receptive field problem.}
A standard convolutional layer with kernel size $k$ and no dilation
has a receptive field of $k$ pixels per layer, growing as $k + (k-1) \cdot (L-1)$
over $L$ layers.
For a sinogram of $N_\theta = 360$ rows, even a very deep standard
CNN would struggle to correlate information across all angles.
We solve this with \textbf{dilated convolutions}, where a layer with
dilation $d$ has kernel spacing $d$, giving an effective kernel
size of $k + (k-1)(d-1)$ per layer, and exponentially growing
receptive field with depth.

\paragraph{Architecture.}
The ConvolutionalDomainTransformer (CDT) first performs a resolution
alignment:
\begin{equation}
    \myvec{x}_{\mathrm{grid}} =
    \mathrm{BilinearInterp}\!\left(\hat{\myvec{s}}_{\mathrm{clean}},\; W \times H\right)
    \label{eq:interp}
\end{equation}
This operation is \textbf{resolution-independent}: the sinogram is
rescaled to match the target image spatial grid regardless of its
original shape, enabling the same trained model to operate on
different detector configurations without retraining.

The subsequent feature extraction follows an encoder-bottleneck-decoder
structure with $C_2 = 64$ base channels:
\begin{align*}
    \myvec{e}_0 &= \mathrm{DoubleConv}_{1 \to C_2}(\myvec{x}_{\mathrm{grid}}) \\
    \myvec{e}_1 &= \mathrm{DoubleConv}_{C_2 \to 2C_2}(\mathrm{MaxPool}_{2}(\myvec{e}_0)) \\
    \myvec{e}_2 &= \mathrm{DoubleConv}_{2C_2 \to 4C_2}(\mathrm{MaxPool}_{2}(\myvec{e}_1)) \\
    \myvec{b} &= \mathrm{Bottleneck}(\myvec{e}_2)
\end{align*}

The \textbf{dilated bottleneck} uses two stacked dilated convolutions:
\begin{align}
    \myvec{b}_1 &= \mathrm{ReLU}\!\left(\mathrm{Conv}_{3 \times 3, d=2}^{4C_2 \to 4C_2}(\myvec{e}_2)\right)
        \label{eq:dil1} \\
    \myvec{b} &= \mathrm{ReLU}\!\left(\mathrm{Conv}_{3 \times 3, d=4}^{4C_2 \to 4C_2}(\myvec{b}_1)\right)
        \label{eq:dil2}
\end{align}
With kernel $k=3$ and dilations $d \in \{2, 4\}$, the effective
receptive field of a single bottleneck unit is:
\begin{equation}
    r_{\mathrm{eff}} = 1 + \sum_{l} (k-1) \cdot d_l
    = 1 + 2 \cdot 2 + 2 \cdot 4 = 13 \;\text{px (per stage)}
\end{equation}
Across three encoder + two bottleneck + two decoder stages with
$2\times$ spatial reductions, the total effective receptive field on
the input sinogram exceeds $256$ pixels in both spatial dimensions
for a $256 \times 256$ target, ensuring full angular coverage.

The decoder uses skip connections identical in spirit to the
classical U-Net:
\begin{align*}
    \myvec{d}_1 &= \mathrm{DoubleConv}_{4C_2 \to 2C_2}\!\left(\mathrm{BilinearUp}(\myvec{b}) + \myvec{e}_1\right) \\
    \myvec{d}_0 &= \mathrm{DoubleConv}_{2C_2 \to C_2}\!\left(\mathrm{BilinearUp}(\myvec{d}_1) + \myvec{e}_0\right) \\
    \myvec{r} &= \mathrm{Conv}_{1 \times 1}^{C_2 \to 1}(\myvec{d}_0)
\end{align*}

The output $\myvec{r} \in \mathbb{R}^{1 \times W \times H}$ is the
\textit{rough image}: a coarse but topologically correct
reconstruction of the CT slice produced entirely by learned
transformations.

\begin{figure}[H]
    \centering
    \figplaceholder{%
        Figure 4: ConvolutionalDomainTransformer (CDT) detailed diagram.
        Show: Clean sinogram (Angles$\times$Det) $\to$ bilinear resize to W$\times$H
        $\to$ DoubleConv $\to$ EncBlock1 $\to$ EncBlock2 $\to$
        Dilated Bottleneck (d=2, then d=4) $\to$ DecBlock1 (skip from Enc1) $\to$
        DecBlock2 (skip from Enc0) $\to$ 1$\times$1 conv $\to$ Rough Image.
        Annotate receptive field growth at each stage.
        [DIAGRAM PENDING]
    }{5cm}
    \caption{
        \textbf{Stage 2: ConvolutionalDomainTransformer.}
        Dilated convolutions in the bottleneck achieve a global
        receptive field, allowing the network to correlate information
        across all angles simultaneously --- the key operation in
        Radon inversion.
        Skip connections from encoder stages preserve fine-grained
        spatial structure lost during downsampling.
    }
    \label{fig:stage2}
\end{figure}

\subsection{Stage 3: ImageUNet --- Image Domain Enhancement}
\label{sec:stage3}

The rough image $\myvec{r}$ produced by the CDT contains correct
large-scale structure but typically lacks high-frequency edge detail,
particularly at material boundaries where the Hounsfield Unit
gradient is steep.
Stage 3 applies a second, image-domain U-Net ($\mathcal{F}_3$) with
$C_3 = 32$ channels, architecturally identical to Stage 1 but
operating on the image domain rather than the sinogram domain.

The residual design is again employed:
\begin{equation}
    \hat{\myvec{x}} = \mathrm{Conv}_{C_3 \to 1}(\myvec{g}_3) + \myvec{r}
\end{equation}
This ensures that the ImageUNet learns to \textit{refine} the rough
image rather than regenerate it, enforcing stability and enabling
faster convergence.
In practice, Stage 3 learns to:
\begin{itemize}[leftmargin=*,noitemsep]
    \item Sharpen material-air and material-material interfaces.
    \item Remove residual checkerboard artefacts from bilinear
    upsampling in Stage 2.
    \item Correct globally distributed low-level intensity bias
    introduced by the domain transform's learned approximation.
\end{itemize}

\subsection{Parameter Count and Complexity}

\begin{table}[H]
    \centering
    \caption{Parameter count per stage ($C_1{=}32$, $C_2{=}64$, $C_3{=}32$).}
    \label{tab:params}
    \begin{tabular}{lrr}
        \toprule
        \textbf{Stage} & \textbf{Module} & \textbf{Parameters} \\
        \midrule
        Stage 1 & SinogramUNet       & $\approx$\,0.8\,M \\
        Stage 2 & Convolutional DT   & $\approx$\,7.2\,M \\
        Stage 3 & ImageUNet          & $\approx$\,0.8\,M \\
        \midrule
        \textbf{Total} &            & $\approx$\,\textbf{8.8\,M} \\
        \bottomrule
    \end{tabular}
\end{table}

% -----------------------------------------------------------------------
\section{Training}
\label{sec:training}
% -----------------------------------------------------------------------

\subsection{Loss Function}

We train the full PureDLPipeline end-to-end with a
\textbf{multi-term L1 loss} that provides gradient signal in both
domains simultaneously:
\begin{align}
    \mathcal{L} &=
        \underbrace{\|\hat{\myvec{s}}_{\text{clean}} - \myvec{s}^{*}\|_1}_{\text{sinogram loss } \mathcal{L}_s}
        + \underbrace{\|\myvec{r} - \myvec{x}^{*}\|_1}_{\text{rough image loss } \mathcal{L}_r}
        + \underbrace{\|\hat{\myvec{x}} - \myvec{x}^{*}\|_1}_{\text{final image loss } \mathcal{L}_f}
    \label{eq:loss}
\end{align}
where $\myvec{s}^{*}$ is the clean full-view sinogram and
$\myvec{x}^{*}$ is the FDK ground-truth reconstruction.

\paragraph{Rationale.}
The $\mathcal{L}_s$ term prevents the sinogram denoiser from
distorting the physics-consistent structure of the measurement.
The $\mathcal{L}_r$ term provides direct gradient feedback to
Stage 2 (the most critical component), ensuring that the bottleneck
receives strong learning signal independent of the Stage 3 output.
The $\mathcal{L}_f$ term drives the final image quality.
L1 is preferred over L2 (MSE) as it is more robust to outliers and
produces sharper edges in image reconstruction tasks~\cite{zhao2017}.

\subsection{Training Dataset: DualDomainDataset}

For each training sample, the dataset stores a triplet:
\begin{equation}
    \{(\myvec{s}_{\text{noisy}}, \myvec{s}^{*}, \myvec{x}^{*})\}_{i=1}^{N}
    \label{eq:triplet}
\end{equation}
When the noisy sinogram has fewer angles than the target
(sparse-view case), it is bilinearly interpolated to the full
angular resolution before being fed to Stage 1.
This allows the sinogram denoiser to learn to \textit{fill in}
missing views as part of its conditioning role.

\subsection{Optimisation}

\begin{table}[H]
    \centering
    \caption{Training hyperparameters.}
    \label{tab:hyper}
    \begin{tabular}{ll}
        \toprule
        \textbf{Hyperparameter} & \textbf{Value} \\
        \midrule
        Optimiser            & Adam \\
        Learning rate        & $10^{-3}$ \\
        Batch size           & 2 (dual-domain, GPU memory bound) \\
        Epochs               & 20 \\
        Validation split     & 20\% (random, seeded) \\
        Validation metric    & L1 loss + PSNR (dB) \\
        Checkpointing        & Best validation-loss epoch \\
        \bottomrule
    \end{tabular}
\end{table}

The model is trained on CUDA when available, with automatic
mixed precision available for future extension.
Best and last checkpoints are saved separately, enabling
resumption of training and model selection based on validation PSNR.

% -----------------------------------------------------------------------
\section{Experimental Setup}
\label{sec:experiments}
% -----------------------------------------------------------------------

\subsection{Dataset}

We generate projections from five 3D STL meshes of industrial objects:
an engineering component (FINAL3, FINAL30), a biological object
(lizard scan), a consumer product (whistle), and a complex mechanical
assembly (Frantic Trug-Crift).
This diversity ensures the network learns geometry-agnostic
features rather than object-specific signatures.

For each object, 360 projections are acquired at $1^\circ$ intervals,
and four degraded subsets are derived:
sparse-view at step 2 (180 views) and step 4 (90 views),
and limited-angle $[0^\circ, 180^\circ]$ and $[30^\circ, 210^\circ]$.
This yields approximately 20 training geometries across 5 objects,
with each geometry producing multiple 2D training slices extracted
from the reconstructed volume.

\subsection{Baselines}

We compare against:
\begin{enumerate}[leftmargin=*, noitemsep]
    \item \textbf{FDK (ASTRA)}: Classical cone-beam FDK reconstruction
    on the degraded sinogram, with a Ram-Lak ramp filter.
    \item \textbf{FDK + ImageUNet}: Classical FDK followed by a
    standalone image-domain U-Net (Stage 3 architecture), trained
    with the same image-domain L1 loss.
    \item \textbf{PureDL (Ours)}: The full three-stage pipeline.
\end{enumerate}

\subsection{Metrics}

We report:
\begin{itemize}[leftmargin=*, noitemsep]
    \item \textbf{PSNR} (Peak Signal-to-Noise Ratio, dB): measures
    pixel-level fidelity to the ground-truth FDK volume.
    \item \textbf{SSIM} (Structural Similarity Index Measure):
    perceptually weighted measure of structural, luminance, and
    contrast similarity~\cite{wang2004ssim}.
\end{itemize}

% -----------------------------------------------------------------------
\section{Results}
\label{sec:results}
% -----------------------------------------------------------------------

\begin{table}[H]
    \centering
    \caption{
        Quantitative comparison on the test set.
        \textbf{[Results pending GPU execution.]}
        Higher PSNR and SSIM are better.
    }
    \label{tab:results}
    \begin{tabular}{lcccc}
        \toprule
        & \multicolumn{2}{c}{\textbf{Sparse-view ($\times$4)}}
        & \multicolumn{2}{c}{\textbf{Limited Angle $[0^\circ,180^\circ]$}} \\
        \cmidrule(lr){2-3}\cmidrule(lr){4-5}
        \textbf{Method} & PSNR$\uparrow$ & SSIM$\uparrow$
                        & PSNR$\uparrow$ & SSIM$\uparrow$ \\
        \midrule
        FDK (ASTRA)            & -- & -- & -- & -- \\
        FDK + ImageUNet        & -- & -- & -- & -- \\
        \textbf{PureDL (Ours)} & -- & -- & -- & -- \\
        \bottomrule
    \end{tabular}
\end{table}

\begin{figure}[H]
    \centering
    \figplaceholder{%
        Figure 5: Qualitative reconstruction comparison.
        Row 1: Sparse-view ($\times4$).
        Row 2: Limited-angle ($[0^\circ, 180^\circ]$).
        Columns (left to right): (a) Input degraded sinogram,
        (b) FDK baseline reconstruction, (c) FDK + ImageUNet,
        (d) PureDL Stage-2 rough image, (e) PureDL final output,
        (f) Ground-truth FDK (full 360 views).
        [IMAGES PENDING GPU EXECUTION]
    }{7cm}
    \caption{
        \textbf{Qualitative Results.}
        Comparison of reconstruction quality across degradation types.
        PureDL is expected to suppress streak artefacts in the
        sparse-view case and reduce the characteristic wedge-shaped
        missing-cone artefacts in the limited-angle case.
    }
    \label{fig:qualitative}
\end{figure}

\begin{figure}[H]
    \centering
    \figplaceholder{%
        Figure 6: Training curves.
        Plot: Train loss and validation loss vs. epoch (20 epochs).
        Secondary axis: Validation PSNR (dB) vs. epoch.
        [PENDING]
    }{4cm}
    \caption{
        \textbf{Training Convergence.}
        Validation PSNR and L1 loss curves over 20 epochs of training
        on the simulated gVXR dataset.
    }
    \label{fig:training}
\end{figure}

\subsection{Ablation Study}

To validate each architectural component, we conduct an ablation:
\begin{itemize}[leftmargin=*, noitemsep]
    \item \textbf{Stage 2 only} (no Stage 1 denoising, no Stage 3 refiner):
    direct sinogram-to-image mapping.
    \item \textbf{Stage 1 + Stage 2} (no Stage 3 refinement).
    \item \textbf{Stage 2 with standard convolutions only} (no dilation
    in bottleneck).
    \item \textbf{Full PureDL} (Stage 1 + Stage 2 dilated + Stage 3).
\end{itemize}
Results (Table~\ref{tab:ablation}) are expected to demonstrate:
(i) Stage 1 denoising improves domain transformer accuracy,
(ii) dilated convolutions are critical for the bottleneck's global
integration capability, and
(iii) Stage 3 provides a measurable PSNR gain, particularly on
high-frequency edges.

\begin{table}[H]
    \centering
    \caption{Ablation study on sparse-view ($\times$4) test set.
    \textbf{[Results pending GPU execution.]}}
    \label{tab:ablation}
    \begin{tabular}{lcc}
        \toprule
        \textbf{Configuration} & PSNR (dB) & SSIM \\
        \midrule
        Stage 2 only (no dilation)       & -- & -- \\
        Stage 2 only (with dilation)     & -- & -- \\
        Stage 1 + Stage 2                & -- & -- \\
        Stage 1 + Stage 2 + Stage 3      & -- & -- \\
        \bottomrule
    \end{tabular}
\end{table}

% -----------------------------------------------------------------------
\section{Discussion}
\label{sec:discussion}
% -----------------------------------------------------------------------

\subsection{On the Learnability of the Radon Inversion}

The classical FBP can be written as a composition of a filtering
operation and a back-projection integral.
The filtering is a fixed ramp function in Fourier space; it is the
back-projection that fundamentally requires knowing the geometry.
Our CDT replaces this two-step process with a single learned
non-linear operator.
The key insight is that the sinogram representation already encodes
all geometric information implicitly: each point in the image space
leaves a sinusoidal trace in the sinogram (Eq.~\ref{eq:radon}).
The network must learn to ``unscramble'' these sinusoidal traces
into a coherent image.
Dilated convolutions enable this by providing receptive fields that
span the full angular dimension of the sinogram, allowing the
network to simultaneously correlate contributions from all angles
--- the operation that back-projection performs analytically.

\subsection{Resolution Independence}

A practically significant property of our CDT is the bilinear
interpolation in Eq.~\ref{eq:interp}: the sinogram is mapped to a
fixed $W \times H$ grid before any learned processing.
This means:
\begin{enumerate}[leftmargin=*, noitemsep]
    \item The model is agnostic to the number of detector pixels
    $N_d$ at inference time.
    \item A model trained on $256 \times 256$ target images can be
    straightforwardly adapted to $512 \times 512$ by simply changing
    \texttt{target\_image\_size}.
    \item Oblique or non-standard angular sampling schemes are
    handled transparently by the initial interpolation.
\end{enumerate}

\subsection{Limitations and Future Work}

\textbf{3D vs.\ 2D.}
The current pipeline operates slice-by-slice on 2D sinograms extracted
from the 3D volume.
Extending to full 3D reconstruction (cone-beam sinogram $\to$ 3D volume)
would require significantly more GPU memory but is architecturally
straightforward by replacing 2D convolutions with 3D variants.

\textbf{Generalisation to unseen geometries.}
The model is trained on a fixed virtual scanner geometry.
At inference on a real scanner, geometric mismatch may reduce
accuracy.
Conditioning the network on scanner parameters (SOD, SDD, pixel
size) as auxiliary inputs is a natural extension.

\textbf{No beam hardening correction.}
The simulation uses monochromatic X-rays.
Real polychromatic beams introduce beam-hardening artefacts
(cupping and streaks near high-density materials) that are not
present in the training data.
Future work should incorporate polychromatic gVXR simulations.

% -----------------------------------------------------------------------
\section{Conclusion}
\label{sec:conclusion}
% -----------------------------------------------------------------------

We have presented the \textbf{PureDLPipeline}, a three-stage deep
learning architecture for CT reconstruction that entirely replaces
the classical Radon-transform inversion (FDK/FBP) with learned
transformations at inference time.
The core novelty is the \textit{ConvolutionalDomainTransformer}:
a resolution-independent, dilated convolutional network whose
bottleneck achieves a global receptive field sufficient to
approximate the angular integration intrinsic to back-projection.

Complementing the architecture is a complete physical simulation
pipeline using gVXR, which generates realistic cone-beam projection
data from CAD meshes with photon-counting noise, providing a
principled and reproducible training environment free of the
acquisition constraints of real scanner data.

The multi-term L1 loss simultaneously optimises sinogram-domain
fidelity (Stage 1), domain transform accuracy (Stage 2), and
final image quality (Stage 3), providing richer gradient information
than image-only losses used in prior work.

Experimental validation on simulated industrial objects across
sparse-view and limited-angle degradation regimes is ongoing;
quantitative and qualitative results will be reported upon GPU
execution completion.
Our open-source implementation is available at:
\begin{center}
\small\url{https://github.com/kitretsu2809/gVXRsimulation_PureDLReconstruction_Pipeline}
\end{center}

% -----------------------------------------------------------------------
\section*{Acknowledgements}
% -----------------------------------------------------------------------

[Acknowledgements to supervisors, funding agencies, or compute
resources to be added.]

% -----------------------------------------------------------------------
\bibliographystyle{plainnat}
\begin{thebibliography}{99}

\bibitem{kak2001}
A.~C. Kak and M.~Slaney,
\textit{Principles of Computerized Tomographic Imaging},
Society for Industrial and Applied Mathematics, 2001.

\bibitem{feldkamp1984}
L.~A. Feldkamp, L.~C. Davis, and J.~W. Kress,
``Practical cone-beam algorithm,''
\textit{Journal of the Optical Society of America A},
vol.~1, no.~6, pp.~612--619, 1984.

\bibitem{jin2017}
K.~H. Jin, M.~T. McCann, E.~Froustey, and M.~Unser,
``Deep convolutional neural network for inverse problems in imaging,''
\textit{IEEE Transactions on Image Processing},
vol.~26, no.~9, pp.~4509--4522, 2017.

\bibitem{chen2017}
H.~Chen, Y.~Zhang, M.~K. Kalra, F.~Lin, Y.~Chen, P.~Liao,
J.~Zhou, and G.~Wang,
``Low-dose CT with a residual encoder-decoder convolutional neural network,''
\textit{IEEE Transactions on Medical Imaging},
vol.~36, no.~12, pp.~2524--2535, 2017.

\bibitem{han2016}
Y.~S. Han and J.~C. Ye,
``Framing U-Net via deep convolutional framelets: Application to sparse-view CT,''
\textit{IEEE Transactions on Medical Imaging},
vol.~37, no.~6, pp.~1418--1429, 2018.

\bibitem{lee2018}
H.~Lee, J.~Lee, H.~Kim, B.~Cho, and S.~Cho,
``Deep-neural-network based sinogram synthesis for sparse-view CT image reconstruction,''
\textit{IEEE Transactions on Radiation and Plasma Medical Sciences},
vol.~3, no.~2, pp.~109--119, 2018.

\bibitem{hu2020}
D.~Hu, J.~Liu, T.~Lv, Q.~Zhao, Y.~Zhang, G.~Quan, J.~Feng,
Y.~Chen, and L.~Luo,
``Hybrid-domain neural network processing for sparse-view CT reconstruction,''
\textit{IEEE Transactions on Radiation and Plasma Medical Sciences},
vol.~5, no.~1, pp.~88--98, 2020.

\bibitem{wu2021drone}
W.~Wu, D.~Hu, C.~Niu, H.~Yu, V.~Vardhanabhuti, and G.~Wang,
``DRONE: Dual-domain residual-based optimization network for
sparse-view CT reconstruction,''
\textit{IEEE Transactions on Medical Imaging},
vol.~40, no.~11, pp.~3002--3014, 2021.

\bibitem{liu2020}
J.~Liu, Y.~Hu, X.~Yang, Y.~Chen, H.~Shu, L.~Luo, C.~Toumoulin,
J.-L. Coatrieux, and L.~Li,
``3D feature constrained reconstruction for low dose CT imaging,''
\textit{IEEE Transactions on Circuits and Systems for Video Technology},
vol.~28, no.~5, pp.~1232--1247, 2018.

\bibitem{vidal2019gvxr}
F.~P. Vidal, J.~M. L\'{e}tang, G.~Peix, and P.~Clackdoyle,
``Assessment of scattered radiation and correction methods in
cone-beam computed tomography using gVirtualXray,''
\textit{Nuclear Instruments and Methods in Physics Research A},
vol.~949, p.~162792, 2020.

\bibitem{astra2016}
W.~van Aarle, W.~J. Palenstijn, J.~Cant, E.~Janssens, F.~Bleichrodt,
A.~Dabravolski, J.~De Beenhouwer, K.~J. Batenburg, and J.~Sijbers,
``Fast and flexible X-ray tomography using the ASTRA toolbox,''
\textit{Optics Express},
vol.~24, no.~22, pp.~25129--25147, 2016.

\bibitem{zhao2017}
H.~Zhao, O.~Gallo, I.~Frosio, and J.~Kautz,
``Loss functions for image restoration with neural networks,''
\textit{IEEE Transactions on Computational Imaging},
vol.~3, no.~1, pp.~47--57, 2017.

\bibitem{wang2004ssim}
Z.~Wang, A.~C. Bovik, H.~R. Sheikh, and E.~P. Simoncelli,
``Image quality assessment: From error visibility to structural similarity,''
\textit{IEEE Transactions on Image Processing},
vol.~13, no.~4, pp.~600--612, 2004.

\end{thebibliography}

\end{document}
